\begin{proposition}[3.1.2]
\label{IV.2.3.1.2-fr}
Sean $A$ un anillo noetheriano, $M$ un $A$-módulo,
$X=\operatorname{Spec}(A)$, $\mathcal F=\widetilde M$~; para que un punto
$x\in X$ esté asociado a $\mathcal F$, es necesario y suficiente que el ideal
primo $\mathfrak j_x$ de $A$ esté asociado al módulo $M$ \textup{(en otras palabras,
sea el \emph{anulador} de un elemento $f\in M$)}.
\end{proposition}

Esto se sigue de la definición (3.1.1) y de Bourbaki, \emph{loc. cit.},
\S~1, núm.~2, cor. de la prop.~5, aplicado a
$S=A\setmin\mathfrak j_x$.

\begin{proposition}[3.1.3]
\label{IV.2.3.1.3-fr}
Sean $X$ un preesquema localmente noetheriano, $\mathcal F$ un
$\mathcal O_X$-módulo cuasicoherente, $x$ un punto de $X$, $Z$ el
subpreesquema cerrado reducido de $X$ cuyo espacio subyacente es
$\overline{\{x\}}$ \textup{(\textbf I, 5.2.1)}. Las condiciones siguientes son equivalentes~:
\begin{enumerate}
\item[\emph{a)}] $x\in\operatorname{Ass}(\mathcal F)$.

\oldpage[IV-2]{37}
\item[\emph{b)}] Existen una vecindad abierta $U$ de $x$ y una sección
$f\in\Gamma(U,\mathcal F)$ tales que $U\cap Z$ sea igual a
$\operatorname{Supp}(\mathcal O_U.f)$.
\item[\emph{b$'$)}] Existen una vecindad abierta $U$ de $x$ y una sección
$f\in\Gamma(U,\mathcal F)$ tales que $U\cap Z$ sea una componente
irreducible de $\operatorname{Supp}(\mathcal O_U.f)$.
\item[\emph{c)}] Existe una vecindad abierta $U$ de $x$ y un submódulo de
$\mathcal F|U$ isomorfo a $\mathcal O_Z|U$ \textup{(identificándose
$\mathcal O_Z$ con un cociente de $\mathcal O_X$)}.
\item[\emph{c$'$)}] Existen una vecindad abierta $U$ de $x$ y un submódulo
coherente $\mathcal G$ de $\mathcal F|U$ tales que $U\cap Z$ sea una
componente irreducible de $\operatorname{Supp}(\mathcal G)$.
\end{enumerate}
\end{proposition}

Es claro que \emph{c)} implica \emph{b)}, pues basta tomar como $f$ el elemento
de $\Gamma(U,\mathcal F)$ que corresponde a la sección unidad de
$\mathcal O_Z|U$. Como $U\cap Z$ es irreducible ($\mathbf{0}_{\mathrm I}$, 2.1.6),
\emph{b)} implica \emph{b$'$)} y \emph{b$'$)} implica \emph{c$'$)}, puesto que
$\mathcal O_X$ es coherente ($\mathbf{0}_{\mathrm I}$, 5.3.4). Para ver que
\emph{c$'$)} implica \emph{a)}, es posible reducirse al caso en que
$U=X=\operatorname{Spec}(A)$ es afín, siendo por tanto $A$ noetheriano, y en
que $\mathcal F=\widetilde M$, donde $M$ es un $A$-módulo, y
$\mathcal G=\widetilde N$, donde $N$ es un submódulo de $M$, de tipo finito.
Se sabe entonces que los elementos minimales de
$\operatorname{Supp}(\mathcal G)$ son los puntos maximales de
$V(\operatorname{Ann}(N))$ ($\mathbf{0}_{\mathrm I}$, 1.7.4), y son también los elementos
minimales de $\operatorname{Ass}(N)$ (Bourbaki, \emph{Alg. comm.}, cap.~IV,
\S~1, núm.~3, cor.~1 de la prop.~7)~; como
$\operatorname{Ass}(N)\subset\operatorname{Ass}(M)
=\operatorname{Ass}(\mathcal F)$, se ve que \emph{c$'$)} implica
\emph{a)}. Por último, \emph{a)} implica \emph{c)} en virtud de (3.1.2),
tomando de nuevo $X$ afín, $\mathcal F=\widetilde M$, y $Z$ definido por el
ideal $\mathfrak j_x$ (con las notaciones de (3.1.2)).

